\documentclass[11pt]{article}
\usepackage[T1]{fontenc}
\usepackage{lmodern}
\usepackage[margin=1in]{geometry}
\usepackage{amsmath,amssymb,amsthm}
\usepackage[colorlinks=true,linkcolor=blue,citecolor=blue,urlcolor=blue]{hyperref}
\hypersetup{pdftitle={Undecidability of unnormalized univariate trace inequalities},pdfauthor={Kaiwen Jin},pdfsubject={An exact undecidability reduction for ordinary pure traces}}
\newtheorem{theorem}{Theorem}
\newtheorem{lemma}[theorem]{Lemma}
\newcommand{\tr}{\operatorname{tr}}
\newcommand{\Nzero}{\mathbb{N}_0}
\title{Undecidability of unnormalized univariate trace inequalities}
\author{Kaiwen Jin\thanks{Email: \href{mailto:yui.ui.siki@gmail.com}{yui.ui.siki@gmail.com}.}}
\date{October 13, 2026}
\begin{document}
\maketitle
\begin{abstract}
We prove that no algorithm decides whether a rational polynomial in the
ordinary traces of powers of one real symmetric matrix is nonnegative for
all finite matrix sizes. The proof encodes nonnegative integers as spectral
multiplicities. Stationary polynomial interpolation controls deviations from
the encoding nodes with an error bound independent of the number of
eigenvalues. A polynomial penalty then transfers an undecidable discrete
sign problem to an unrestricted pure trace inequality. The construction uses
only positive powers and never uses the matrix dimension as a trace variable.
\end{abstract}

\section{The decision problem}

For $d\geq1$ and $F\in\mathbb{Q}[z_1,\ldots,z_d]$, consider the statement
\begin{equation}\label{eq:target}
 F\bigl(\tr A,\tr(A^2),\ldots,\tr(A^d)\bigr)\geq0
 \quad\text{for all }n\geq1\text{ and real symmetric }A\in\mathbb{R}^{n\times n}.
\end{equation}
The trace is ordinary and unnormalized. The polynomial is pure trace: its
value is a scalar, with no additional matrix factor, and its scalar constant
$1$ is not $\tr I=n$. By diagonalization, \eqref{eq:target} is equivalent to
\begin{equation}\label{eq:spectrum}
 F\bigl(p_1(x),\ldots,p_d(x)\bigr)\geq0
 \quad\text{for all }n\geq1\text{ and }x\in\mathbb{R}^n,
 \qquad p_r(x)=\sum_{a=1}^n x_a^r.
\end{equation}

Acevedo, Blekherman, Debus and Riener prove undecidability with several
matrix variables and decidability for normalized traces
\cite[Theorems 6.2 and 6.5]{ABDR}. Their Section 7 leaves the ordinary
univariate case open. The positivity results of Klep, Pascoe and Vol\v{c}i\v{c}
use normalized traces \cite[Introduction, arXiv version v2]{KPV}. The distinction is relevant here:
ordinary traces retain integer multiplicities, whereas normalization divides
those multiplicities by the dimension.

\begin{theorem}\label{thm:main}
There is no algorithm that, on every input consisting of a positive integer
$d$ and a polynomial $F\in\mathbb{Q}[z_1,\ldots,z_d]$, always terminates and
correctly decides \eqref{eq:target}.
\end{theorem}

We prove the theorem by reducing nonnegativity on integer tuples to
\eqref{eq:spectrum}. A spectrum supported on $\{0,1,\ldots,k\}$ can record
$k$ arbitrary integer multiplicities. The main point is to control every
spectrum away from these nodes, uniformly in its size. Interpolating functions
with zero first derivatives at the nodes give the required control. Including
$0$ as a node makes the relevant polynomial sums independent of zero padding
and avoids any use of a zeroth power sum.

\section{A uniform estimate for spectral multiplicities}

Write $\Nzero=\{0,1,2,\ldots\}$. Fix $k\geq1$. For each real $x$, let
$j(x)$ be a closest element of $\{0,\ldots,k\}$, choosing the smaller index at
a tie. For a rational polynomial $q(x)=\sum_r b_r x^r$ and $R\geq1$, use the
coefficient bound $\|q\|_R=\sum_r |b_r|R^r$.

\begin{lemma}\label{lem:rounding}
One can effectively construct rational polynomials $e_1,\ldots,e_k,h$ and a
positive rational number $C$, depending only on $k$, such that
\[
 \deg e_i\leq2k+1,\qquad \deg h=2k+2,\qquad
 e_i(0)=h(0)=0,\qquad h(x)\geq0\quad(x\in\mathbb{R}),
\]
and
\begin{equation}\label{eq:pointwise}
 |e_i(x)-\delta_{i,j(x)}|\leq C h(x)
 \quad (x\in\mathbb{R},\ 1\leq i\leq k).
\end{equation}
Consequently, for every $n\geq1$ and $x\in\mathbb{R}^n$, the quantities
\begin{equation}\label{eq:counts}
 s_i=\sum_{a=1}^n e_i(x_a),\qquad
 t=\sum_{a=1}^n h(x_a),\qquad
 m_i=\#\{a:j(x_a)=i\}
\end{equation}
satisfy $m\in\Nzero^k$, $t\geq0$ and $|s_i-m_i|\leq Ct$.
Every $m\in\Nzero^k$ is realized with $s=m$ and $t=0$ by a finite spectrum.
\end{lemma}

\begin{proof}
For $1\leq i\leq k$, define
\begin{equation}\label{eq:interpolants}
 \ell_i(x)=\prod_{\substack{0\leq j\leq k\\j\ne i}}\frac{x-j}{i-j},
 \qquad
 e_i(x)=\bigl[1-2\ell_i'(i)(x-i)\bigr]\ell_i(x)^2,
 \qquad
 h(x)=\prod_{j=0}^k(x-j)^2.
\end{equation}
The squared Lagrange factors give
$e_i(j)=\delta_{ij}$ and $e_i'(j)=0$ at every $j\ne i$. At $i$, the two
terms in the derivative cancel, so $e_i(i)=1$ and $e_i'(i)=0$ as well.
The degree and zero-constant assertions follow from these formulas.

For every $i,j$ the quotient
\[
 q_{ij}(x)=\frac{e_i(x)-\delta_{ij}}{(x-j)^2}
\]
is a rational polynomial of degree at most $2k-1$. Put $R=2(k+1)$ and take
\begin{equation}\label{eq:C}
 C=1+2^{2k}\max_{i,j}\|q_{ij}\|_R
      +2^{2k+2}\max_i\bigl(1+\|e_i\|_1\bigr),
\end{equation}
where $1\leq i\leq k$ and $0\leq j\leq k$. This is a computable rational
number.

If $|x|\leq R$ and $j=j(x)$, every other node $l$ satisfies $|x-l|\geq1/2$.
Indeed, a node at distance less than $1/2$ would be the unique closest node.
It follows that
\[
 h(x)\geq2^{-2k}(x-j)^2,
 \qquad
 |e_i(x)-\delta_{ij}|\leq(x-j)^2\|q_{ij}\|_R\leq C h(x).
\]
This also holds directly when $x=j$. If $|x|>R$, then
$|x-l|\geq |x|-k>|x|/2$ for every node $l$. Therefore
\[
 h(x)\geq (|x|/2)^{2k+2},\qquad
 |e_i(x)-\delta_{i,j(x)}|\leq(1+\|e_i\|_1)|x|^{2k+1},
\]
which again proves \eqref{eq:pointwise}. Summing the pointwise bounds gives
$|s_i-m_i|\leq Ct$ without a factor depending on $n$.
Finally, take one zero entry and $m_i$ copies of $i$ for each $i=1,\ldots,k$.
This spectrum has size $1+\sum_i m_i\geq1$ and gives $s=m$, $t=0$.
\end{proof}

\section{The effective reduction}

The following discrete decision problem is undecidable: given an integer
polynomial, determine whether it is nonnegative on every tuple of positive
integers \cite[Section 5, p.~11 of the arXiv version]{HN}. This is a consequence
of Matiyasevich's theorem used in that source. Replacing each variable by
$u_i+1$ shows that nonnegativity on $\Nzero^k$, for arbitrary finite $k$, is
also undecidable.

We first control the change in such a polynomial when the interpolated
sums $s$ are replaced by the associated integer counts $m$ from
\eqref{eq:counts}.

\begin{lemma}\label{lem:variation}
Let $P(u)=\sum_\alpha c_\alpha u^\alpha\in\mathbb{Z}[u_1,\ldots,u_k]$,
where $|\alpha|=\sum_i\alpha_i$. Set
\[
 D=\max\{1,\deg P\},\qquad
 K=\sum_\alpha |c_\alpha|\,|\alpha|,\qquad
 L=3^{D-1}C^2K^2,
\]
with $D=1$ when $P=0$. If $t\geq0$ and $|s_i-m_i|\leq Ct$ for every $i$,
then
\begin{equation}\label{eq:variation}
 \bigl(P(s)-P(m)\bigr)^2
 \leq L t^2\left(1+\sum_{i=1}^k s_i^2+C^2t^2\right)^{D-1}.
\end{equation}
\end{lemma}

\begin{proof}
For a nonconstant monomial, telescope the difference of its $|\alpha|$
coordinate factors, including repetitions. Each changed factor contributes
at most $Ct$, and the other factors have magnitude at most
$\max_i|s_i|+Ct$. Thus
\[
 |s^\alpha-m^\alpha|
 \leq |\alpha|Ct\bigl(1+\max_i|s_i|+Ct\bigr)^{|\alpha|-1}.
\]
The base is at least $1$. Multiplying by $|c_\alpha|$, summing, and using
$|\alpha|\leq D$ gives
\[
 |P(s)-P(m)|\leq CKt\bigl(1+\max_i|s_i|+Ct\bigr)^{D-1}.
\]
The inequality $(1+a+b)^2\leq3(1+a^2+b^2)$ and
$(\max_i|s_i|)^2\leq\sum_i s_i^2$ prove \eqref{eq:variation}.
For a constant polynomial, including $P=0$, both $K$ and the difference are
zero, and the same bound holds with $D=1$.
\end{proof}

\begin{proof}[Proof of Theorem~\ref{thm:main}]
Given $P\in\mathbb{Z}[u_1,\ldots,u_k]$, set $d=2k+2$ and take the
polynomials and $C$ of Lemma~\ref{lem:rounding}. Write
\[
 e_i(x)=\sum_{r=1}^d e_{ir}x^r,\qquad
 h(x)=\sum_{r=1}^d h_r x^r,
\]
allowing zero coefficients. Define linear forms in the positive power-sum
variables by
\[
 S_i(z)=\sum_{r=1}^d e_{ir}z_r,\qquad
 T(z)=\sum_{r=1}^d h_r z_r.
\]
With $D,K,L$ as in Lemma~\ref{lem:variation}, form the rational polynomial
\begin{equation}\label{eq:F}
 \begin{split}
 F_P(z)={}&2P\bigl(S_1(z),\ldots,S_k(z)\bigr)+1\\
 &+L T(z)^2\left(1+\sum_{i=1}^k S_i(z)^2+C^2T(z)^2\right)^{D-1}.
 \end{split}
\end{equation}
All these operations terminate in exact rational polynomial arithmetic:
finite products, derivatives, exact divisions by $(x-j)^2$, rational
coefficient bounds and finite compositions. The expression can be expanded
into a finite coefficient list. No discrete solution is needed to construct
it.

At $z=(p_1(x),\ldots,p_d(x))$, the forms $S_i$ and $T$ equal the spectral
sums $s_i$ and $t$ in \eqref{eq:counts}. Since $e_i(0)=h(0)=0$, these sums
use only positive powers. The constants in \eqref{eq:F} are scalar
polynomial constants, so the construction never introduces $\tr I$ or a
zeroth power-sum variable.

Suppose $P(m)\geq0$ for every $m\in\Nzero^k$. For any finite real spectrum,
let $m$ be its rounded counts and put $\eta=P(s)-P(m)$. Lemma~\ref{lem:variation}
gives
\[
 \begin{split}
 F_P\bigl(p_1(x),\ldots,p_d(x)\bigr)
 &\geq 2P(m)+2\eta+1+\eta^2\\
 &=2P(m)+(\eta+1)^2\geq0.
 \end{split}
\]
Conversely, if $P(m)<0$ at some $m\in\Nzero^k$, then $P(m)\leq-1$ by
integrality. The node spectrum from Lemma~\ref{lem:rounding} has $s=m$,
$t=0$ and therefore
\[
 F_P\bigl(p_1(x),\ldots,p_d(x)\bigr)=2P(m)+1\leq-1.
\]
We have proved
\[
 P\geq0\text{ on }\Nzero^k
 \quad\Longleftrightarrow\quad
 F_P\bigl(p_1(x),\ldots,p_d(x)\bigr)\geq0
 \text{ for every }n\geq1,\ x\in\mathbb{R}^n.
\]
An always-terminating decider for \eqref{eq:target}, applied to the computable
input $(2k+2,F_P)$, would decide the discrete problem. This contradiction
proves the theorem.
\end{proof}

\paragraph{AI-use statement.}
OpenAI Codex was used to develop the argument, check cited sources and draft
this manuscript. Independent human mathematical review is requested.

\begin{thebibliography}{9}
\bibitem{ABDR}
J.~Acevedo, G.~Blekherman, S.~Debus and C.~Riener,
\emph{The wonderful geometry of the Vandermonde map},
Foundations of Computational Mathematics \textbf{26} (2026), 2005--2051.
\href{https://doi.org/10.1007/s10208-025-09718-6}{doi:10.1007/s10208-025-09718-6}.
\bibitem{KPV}
I.~Klep, J.~E.~Pascoe and J.~Vol\v{c}i\v{c},
\emph{Positive univariate trace polynomials},
Journal of Algebra \textbf{579} (2021), 303--317.
\href{https://doi.org/10.1016/j.jalgebra.2021.03.027}{doi:10.1016/j.jalgebra.2021.03.027}.
The cited version is \href{https://arxiv.org/abs/2009.10081v2}{arXiv:2009.10081v2}.
\bibitem{HN}
H.~Hatami and S.~Norine,
\emph{Undecidability of linear inequalities in graph homomorphism densities},
arXiv:1005.2382v3 (2010).
\url{https://arxiv.org/abs/1005.2382v3}.
\end{thebibliography}
\end{document}
